% Problem record op_89fb664ba06ba5de. This TeX file is derived from
% database/problems_json/op_89fb664ba06ba5de.json by scripts/sync-tex.mjs; edit the JSON record
% and rerun the script rather than editing this file.
\section{Candidate pure-loss second-order converse}
\label{sec:op_89fb664ba06ba5de}

\paragraph{Problem.}
Does the pure-loss bosonic channel admit the following candidate second-order
classical converse under a maximum-photon-number occupation constraint?  Let
$\mathcal N_\eta$ be the single-mode pure-loss channel with transmissivity
$0<\eta<1$, defined in the Heisenberg picture by
\begin{equation}
  \hat b=\sqrt{\eta}\,\hat a+\sqrt{1-\eta}\,\hat e,
  \label{eq:p76-pure-loss-channel}
\end{equation}
where the environment mode $E$ is in the vacuum.  In an $n$-use code, the
channel in Eq.~\eqref{eq:p76-pure-loss-channel} is used to transmit one of
$M$ input states $\rho_m^{A^n}$, with a decoding POVM
$\{\Lambda_m^{B^n}\}_{m=1}^M$.  Write
$\overline\rho_{A^n}:=M^{-1}\sum_m\rho_m^{A^n}$ and let
$\Pi_{\lceil nN_S\rceil}$ project onto the $n$-mode subspace of total photon
number at most $\lceil nN_S\rceil$.  For fixed $N_S>0$,
$\varepsilon\in(0,1)$, and $c>0$, impose
\begin{equation}
  \frac1M\sum_{m=1}^M
    \operatorname{Tr}\!\left[
      \Lambda_m\mathcal N_\eta^{\otimes n}(\rho_m)
    \right]
    \geq1-\varepsilon,
  \qquad
  \operatorname{Tr}\!\left[
    \Pi_{\lceil nN_S\rceil}\overline\rho_{A^n}
  \right]
    \geq1-\delta_n,
  \qquad
  0\leq\delta_n\leq2^{-cn}.
  \label{eq:p76-code-constraints}
\end{equation}
Let $M^*_{\rm occ}(n,\eta,N_S,\varepsilon,c)$ be the largest $M$ satisfying
Eq.~\eqref{eq:p76-code-constraints}.  Define the thermal entropy and its
entropy variance by
\begin{equation}
  g(x):=(x+1)\log_2(x+1)-x\log_2x,
  \qquad
  v(x):=x(x+1)
    \left[\log_2(x+1)-\log_2x\right]^2,
  \label{eq:p76-thermal-entropy-and-variance}
\end{equation}
where $0\log_2 0:=0$.  With the functions in
Eq.~\eqref{eq:p76-thermal-entropy-and-variance}, is the following upper bound
valid as $n\to\infty$?
\begin{equation}
  \log_2 M^*_{\rm occ}(n,\eta,N_S,\varepsilon,c)
  \leq ng(\eta N_S)
    +\sqrt{n\,v(\eta N_S)}\,\Phi^{-1}(\varepsilon)
    +O(\log n),
  \label{eq:p76-second-order-converse}
\end{equation}
Here $\Phi^{-1}$ in Eq.~\eqref{eq:p76-second-order-converse} is the inverse
standard-normal cumulative distribution function, and the implicit constant
may depend on $\eta,N_S,\varepsilon,$ and $c$, but not on $n$.

\subsection*{Status}

Solved

\subsection*{Source}

Wilde, Renes, and Guha leave open whether their Gaussian second-order
expression is an upper bound and observe that such a converse should impose a
photon-number occupation constraint similar to Wilde and Winter's.
Equation~\eqref{eq:p76-code-constraints} adopts a particular exponentially
strong version of that suggestion, so Eq.~\eqref{eq:p76-second-order-converse}
is a precise strengthened formulation rather than a verbatim conjecture from
their paper
\sourcecite{ref:p76-wilde-renes-guha}{WRG16}.

\subsection*{Progress}

\begin{itemize}
  \item Yue Tu and Yifan Jing showed that Eq.~\eqref{eq:p76-second-order-converse}
fails for every $0<\eta<1$, $N_S>0$, $c>0$, and $0<\varepsilon<\tfrac12$.
Each codeword carries exactly $L=\lceil nN_S\rceil$ photons in a condensate
$\lvert u;L\rangle=(\sum_{j=1}^n u_j\hat a_j^\dagger)^L\lvert0\rangle/\sqrt{L!}$
along a unit direction $u\in\mathbb C^n$, so the occupation constraint in
Eq.~\eqref{eq:p76-code-constraints} holds with $\delta_n=0$.  Pure loss maps
$\lvert u;L\rangle$ to a binomial mixture of the pure condensates
$\lvert u;k\rangle$, and $\langle u;k\vert w;k\rangle=\langle u\vert w\rangle^k$.
Random directions, photon counting, and a square-root measurement in each
photon-number sector give
\begin{equation}
  \log_2 M^*_{\rm occ}(n,\eta,N_S,\varepsilon,c)
  \geq ng(\eta N_S)+\sqrt{n V_{\rm occ}}\,\Phi^{-1}(\varepsilon)-O(\log n),
  \qquad
  V_{\rm occ}=\frac{1-\eta}{1+\eta N_S}\,v(\eta N_S).
  \label{eq:p76-fixed-photon-achievability}
\end{equation}
Since $V_{\rm occ}<v(\eta N_S)$ and $\Phi^{-1}(\varepsilon)<0$ for
$\varepsilon<\tfrac12$, Eq.~\eqref{eq:p76-fixed-photon-achievability} exceeds
the right side of Eq.~\eqref{eq:p76-second-order-converse} by a positive
multiple of $\sqrt n$.  A photon-counting projector bound gives the matching
converse, so under this constraint the dispersion is $V_{\rm occ}$ rather
than $v(\eta N_S)$.  The disproof is formally verified in Lean~4 with
Mathlib; its only modeling input is the standard Fock-basis Kraus form of the
pure-loss channel.  The matching converse and the $O(\log n)$ remainder are
proved on paper only, in the
\href{https://github.com/Naixu-Guo/quantum-open-problems/blob/main/database/proofs/pure-loss-second-order.md}{proof note}
\sourcecite{ref:p76-tu-jing}{TJ26}.

  \item Wilde, Renes, and Guha proved the corresponding second-order
  achievability expression
  \begin{equation}
    \log_2 M^*(\mathcal N_\eta^{\otimes n},N_S,\varepsilon)
    \geq ng(\eta N_S)
      +\sqrt{n\,v(\eta N_S)}\,\Phi^{-1}(\varepsilon)
      +O(\log n).
    \label{eq:p76-second-order-achievability}
  \end{equation}
  Their basic derivation of
  Eq.~\eqref{eq:p76-second-order-achievability} uses a mean-photon-number
  ensemble.  They also modify the construction to obtain exponentially small
  leakage outside the cutoff subspace, but with degraded second-order
  parameters; neither argument supplies
  Eq.~\eqref{eq:p76-second-order-converse}
  \sourcecite{ref:p76-wilde-renes-guha}{WRG16}.

  \item Wilde and Winter proved a first-order strong converse at the rate
  $g(\eta N_S)$ under the maximum-photon-number occupation constraint.  For
  every fixed energy and rate backoff, their modified coherent-state codes
  have exponentially small cutoff leakage.  This does not establish the
  exact $g(\eta N_S)$ first-order term and the stated dispersion while also
  maintaining an arbitrary prescribed exponent $c$ in
  Eq.~\eqref{eq:p76-code-constraints}.  Moreover, a first-order strong
  converse controls rates a fixed distance above capacity and does not
  determine the
  $\sqrt n$ coefficient in Eq.~\eqref{eq:p76-second-order-converse}
  \sourcecite{ref:p76-wilde-winter}{WW14}.

  \item Under only a mean-photon-number constraint, Wilde and Winter
  constructed codes with nonvanishing success probability at rates above
  $g(\eta N_S)$.  Hence a converse of the form
  Eq.~\eqref{eq:p76-second-order-converse} is false under that weaker
  constraint; the occupation hypothesis in
  Eq.~\eqref{eq:p76-code-constraints} is essential
  \sourcecite{ref:p76-wilde-winter}{WW14}.

  \item Historical GitHub report (2026-09-21): Yuxuan Zhang (GitHub: yuxuanzhang1995) reported a fixed-photon-number coding construction claimed to contradict the proposed pure-loss second-order converse. \href{https://github.com/Naixu-Guo/quantum-open-problems/issues/91}{Issue \#91}.

  \item Historical GitHub report (2026-09-23): Ivann1242 reported Yue Tu and Yifan Jing's fixed-photon condensate construction and proposed pure-loss dispersion result. \href{https://github.com/Naixu-Guo/quantum-open-problems/pull/99}{Pull request \#99}.
\end{itemize}

\subsection*{References}

\begin{enumerate}[label={},leftmargin=4.8em,labelsep=0.5em]
  \footnotesize
  \item[\textup{[WRG16]}]\label{ref:p76-wilde-renes-guha}
  M. M. Wilde, J. M. Renes, and S. Guha, ``Second-Order Coding Rates for
  Pure-Loss Bosonic Channels,'' \emph{Quantum Information Processing}
  \textbf{15}, 1289--1308 (2016).
  \href{https://doi.org/10.1007/s11128-015-0997-x}{doi:10.1007/s11128-015-0997-x};
  \href{https://arxiv.org/abs/1408.5328}{arXiv:1408.5328}.

  \item[\textup{[WW14]}]\label{ref:p76-wilde-winter}
  M. M. Wilde and A. Winter, ``Strong Converse for the Classical Capacity of
  the Pure-Loss Bosonic Channel,'' \emph{Problems of Information
  Transmission} \textbf{50}, 117--132 (2014).
  \href{https://doi.org/10.1134/S003294601402001X}{doi:10.1134/S003294601402001X};
  \href{https://arxiv.org/abs/1308.6732}{arXiv:1308.6732}.

  \item[\textup{[TJ26]}]\label{ref:p76-tu-jing}
  Y. Tu and Y. Jing, ``Fixed-photon condensate codes and the pure-loss
second-order converse,'' note for QIQCOP Zoo (23 September 2026), with a
Lean~4 formalization.
\href{https://github.com/Naixu-Guo/quantum-open-problems/blob/main/database/proofs/pure-loss-second-order.md}{Proof note};
\href{https://github.com/Ivann1242/pure-loss-second-order-lean}{Lean verification}.
\end{enumerate}

\subsection*{Comment}

Equation~\eqref{eq:p76-second-order-converse} is settled negatively for every
$0<\varepsilon<\tfrac12$ by the fixed-photon codes above.  For
$\varepsilon>\tfrac12$ the stated bound is valid but not sharp, and at
$\varepsilon=\tfrac12$ the second-order terms vanish; both statements rely on
the paper converse.  The mean-photon-number ensemble of
\sourcecite{ref:p76-wilde-renes-guha}{WRG16} is therefore not second-order
optimal under the occupation constraint.  No external peer review or
historical-priority claim is asserted
\sourcecite{ref:p76-tu-jing}{TJ26}.

\subsection*{Field}

Quantum Communication

\subsection*{Topic}

Bosonic channels; Classical capacity; One-shot and finite-blocklength bounds

\subsection*{ID}

\texttt{op\_89fb664ba06ba5de}
